\section{Assumptions, limitations and open items}
\label{sec:assumptions}

\subsection{Assumptions}
\begin{table}[htbp]
  \centering
  \small
  \renewcommand{\arraystretch}{1.25}
  \begin{tabularx}{\linewidth}{@{}l X@{}}
    \toprule
    \textbf{Assumption} & \textbf{Consequence and justification} \\
    \midrule
    Purely resistive DC line
      & No inductance or capacitance appears in the equations. Correct for a
        steady-state DC flow; line charging and inductance matter only for
        transients and fault studies. \\
    Constant conductor resistance
      & The specific resistance is taken at the temperature of the conductor
        table. Resistance rises roughly 0.4\,\%/K for aluminium, so a hot-day
        loss figure is a few percent higher than reported here. Handled by
        choosing conductor data at the design temperature. \\
    Ideal voltage regulation at the rectifier
      & $V_{d}$ holds exactly at the sending terminals irrespective of
        loading. Real control holds it within a droop band. \\
    Station losses proportional to throughput
      & No no-load term, so light-load efficiency is overstated and full-load
        efficiency is representative. Adequate for screening; a two-term
        model ($p_{0} + p_{1}P$) is the natural refinement. \\
    Balanced bipole carries no DMR current
      & Pole currents are taken as exactly equal. In service a small
        unbalance flows in the DMR; its loss contribution is second order
        (the square of a few percent of $I_{d}$). \\
    One pole out at unchanged pole current
      & Keeps the outage case inside the ratings established for normal
        operation. An alternative case at unchanged \emph{power} would double
        the pole current and require its own rating basis. \\
    Single operating point
      & The tab evaluates one dispatch, not a duty cycle. Annual energy losses
        require a loading profile and a separate loss-energy calculation. \\
    \bottomrule
  \end{tabularx}
  \caption{Assumptions underlying the load flow calculation.}
  \label{tab:assumptions}
\end{table}

\subsection{Outside the scope}
The calculation deliberately stops short of the following, each of which
belongs to a study rather than to a spreadsheet cell:

\begin{itemize}\setlength{\itemsep}{0.3em}
  \item \textbf{Reactive power and AC-side voltage.} Only active power crosses
        the converter boundary in this model. Converter reactive capability and
        AC network voltage support are covered by the AC calculations on the
        Power tab and, ultimately, by interconnection studies.
  \item \textbf{Ground and electrode return.} Only metallic return is
        modelled. A bipole operating with earth return has a different loop
        resistance and a set of electrode, corrosion and interference
        constraints that sit outside a load flow.
  \item \textbf{Thermal rating of the conductor.} The tab reports the current
        that flows; whether the conductor can carry it continuously at the
        ambient and wind conditions of the route is an ampacity calculation.
  \item \textbf{Converter overload and short-term ratings.} The margin of
        Equation~\eqref{eq:margin} is a continuous rating margin, not an
        overload capability.
  \item \textbf{Multi-terminal and intermediate taps.} The equations assume a
        two-terminal, point-to-point link.
  \item \textbf{Dynamics.} Control response, ramp rates, pole-block transients
        and the transition into the one-pole-out condition are all outside a
        steady-state flow.
\end{itemize}

\subsection{Open items}
\begin{enumerate}[label=(\arabic*)]
  \item \textbf{DMR conductor selection.} The reference case uses a single
        Bluebird conductor as a placeholder and the outage-case loss is
        sensitive to it (Section~\ref{sec:reference}). To be replaced with the
        selected DMR conductor and bundle.
  \item \textbf{Unit label on the Losses tab.} The specific resistance row is
        labelled $\Omega$/mi but holds $\Omega$/1000\,ft; the $5.28$ factor in
        the loss formula confirms it. The label should be corrected so the two
        tabs read consistently.
  \item \textbf{Converter loss figure.} 0.8\,\% per station is a generic
        VSC-class value. To be replaced with vendor loss curves, at which
        point the two-term loss model of Table~\ref{tab:assumptions} becomes
        worthwhile.
  \item \textbf{Temperature basis for conductor resistance.} To be stated
        explicitly alongside the conductor data, so the reported losses carry
        a defined temperature.
  \item \textbf{Loss-energy evaluation.} A loading profile and an annual
        loss-energy figure are the natural extension, since it is energy
        rather than peak loss that enters an economic comparison of voltage
        levels and conductor sizes.
\end{enumerate}
